\chapter{MultiNest: Nested Sampling from the Ground Up}
\label{ch:multinest}

\newthought{This chapter is for anyone who needs a Bayesian evidence}, not just a map
of the parameter space --- model comparison, a significance from an evidence ratio, or
a properly weighted posterior. It is also the chapter that teaches nested sampling
itself: the algorithm's core idea, exactly how \Jtwo's code implements and distributes
it, and how to configure it, using \sampler{MultiNest} --- the simpler, fixed-budget
member of the family --- as the vehicle. Chapter~\ref{ch:dynesty} builds directly on
everything here and covers only what its dynamic engine, \sampler{Dynesty}, adds.
After this chapter you can explain what a nested-sampling run is actually computing,
read its \textsc{yaml} block fluently, and recognise every line it prints while
running.

\section{What nested sampling computes}
\label{sec:nested-what}

A Bayesian evidence is a $d$-dimensional integral over your whole parameter space:
$Z = \int L(\theta)\,\pi(\theta)\,d\theta$, where $\pi$ is the prior and $L$ the
likelihood. Nested sampling's insight, due to Skilling\cite{Skilling:2006gxv}, is to
collapse that hard $d$-dimensional integral into an easy \emph{one}-dimensional one.
Define the \emph{enclosed prior volume} at likelihood threshold $\lambda$ as
$X(\lambda) = \int_{L(\theta) > \lambda} \pi(\theta)\,d\theta$ --- the fraction of the
prior with likelihood above $\lambda$. As $\lambda$ rises from $0$ to $L_{\max}$, $X$
falls monotonically from $1$ to $0$. Inverting this relation to $L(X)$ turns the
evidence into
\[
Z = \int_0^1 L(X)\,dX,
\]
a single integral over the unit interval that nested sampling estimates numerically ---
by literally constructing a decreasing sequence of $X$ values with known likelihoods.
Figure~\ref{fig:nested-loop} is the loop that constructs them.

\begin{marginfigure}
\centering
\smallcaps{$N$ live points}\\[2pt]
$\big\downarrow$ \emph{retire the lowest $L$}\\[2pt]
\smallcaps{record it; $X$ shrinks}\\[2pt]
$\big\downarrow$ \emph{draw a replacement with $L$ higher}\\[2pt]
\smallcaps{repeat until converged}\\[2pt]
$\big\downarrow$\\[2pt]
\smallcaps{evidence $Z$ + posterior}
\caption{The nested-sampling loop. Every retired point is a slice of shrinking prior
volume at a known likelihood --- exactly the data $Z=\int_0^1 L(X)\,dX$ needs.}
\label{fig:nested-loop}
\end{marginfigure}

\subsection{The algorithm, step by step}
\label{sec:nested-algorithm}

\begin{enumerate}
  \item \textbf{Initialise.} Draw $N$ \emph{live points} (\yamlkey{nlive}) from the
        prior $\pi$. Their enclosed volume starts at $X_0 = 1$.
  \item \textbf{Retire the worst.} Find the live point with the lowest likelihood
        $L_i$; record it as a \emph{dead point} together with its enclosed volume.
        Each retirement shrinks the enclosed volume by an expected factor of
        $N/(N+1)$, so after $i$ retirements $X_i \approx e^{-i/N}$.
  \item \textbf{Replace it.} Draw a fresh point from the prior, but
        \emph{constrained to $L(\theta) > L_i$} --- it must land strictly inside the
        shrinking high-likelihood region. This constrained draw is the one genuinely
        hard step, and it is where \sampler{MultiNest}'s contribution lives
        (\S\ref{sec:multinest-ellipsoids}).
  \item \textbf{Accumulate the evidence.} Each dead point contributes
        $L_i \cdot w_i$ to $Z$, where $w_i \approx X_{i-1} - X_i$ is the volume slice
        it represents. The running sum is exactly a numerical quadrature of
        $\int_0^1 L(X)\,dX$.
  \item \textbf{Stop.} The live points still sitting above the current threshold could
        still contribute up to $L_{\max}\cdot X_{\text{current}}$ more evidence. Stop
        once that \emph{possible} remaining contribution is smaller than
        \yamlkey{dlogz} relative to the evidence already banked --- and fold the
        surviving live points into the evidence and the posterior
        (\yamlkey{add\_live: true}).
\end{enumerate}

Every dead \emph{and} live point doubles as a posterior sample: its importance weight
is exactly $L_i w_i / Z$. One run produces both numbers you came for.

\subsection{MultiNest's contribution: bounding by ellipsoids}
\label{sec:multinest-ellipsoids}

Step 3 above --- draw a point with $L(\theta) > L_i$ --- is where a naive
implementation dies. Rejecting from the \emph{full} prior until one lands above the
threshold works early on, but the high-likelihood region shrinks exponentially in
$X$, so the acceptance rate collapses to zero as the run progresses.

\sampler{MultiNest}, introduced by Feroz, Hobson \& Bridges\cite{Feroz:2008xx}, solves
this by \textbf{bounding the live points with a union of ellipsoids} instead of
sampling from the raw prior:

\begin{enumerate}
  \item Cluster the current live points (splitting automatically around separate
        modes when the posterior is multimodal).
  \item Fit one ellipsoid per cluster, enlarged by a safety margin
        (\yamlkey{enlarge}) so the true constrained region is never excluded.
  \item Draw new candidates \emph{uniformly from the union of these ellipsoids}
        (correcting for their overlap), and accept the first one with
        $L(\theta) > L_i$.
\end{enumerate}

Because the ellipsoids track the shrinking high-likelihood region far more tightly
than the full prior does, the acceptance rate stays usable for the whole run --- and
because there can be \emph{several} ellipsoids, the same mechanism handles multimodal
posteriors without extra bookkeeping. This is literally what \yamlkey{bound: multi}
selects in the \textsc{yaml} block below; it is also the default in every card in this
book.

\begin{jnote}
The \yamlkey{bound}/\yamlkey{sample} vocabulary you configure in
\S\ref{sec:nested-sampler-block} names exactly these choices: \yamlkey{bound: multi}
is Feroz \& Hobson's union-of-ellipsoids region; \yamlkey{sample: unif} is uniform
rejection sampling inside it --- \sampler{MultiNest}'s original proposal mechanism.
\yamlkey{sample: rwalk} and the \yamlkey{slice} family are newer proposal mechanisms
(a short random walk, or slice sampling, confined to the bound) that stay efficient in
higher dimensions, where naive rejection inside even a well-fitted ellipsoid starts to
struggle again.
\end{jnote}

\section{How the code implements it}
\label{sec:multinest-code}

\Jtwo\ does not reimplement nested sampling --- it takes the reference
\code{dynesty} package (whose author's own paper, \cite{Speagle:2019ivv}, is required
reading for Chapter~\ref{ch:dynesty}'s dynamic engine) and threads it onto the
distributed runtime. Understanding that wiring explains every YAML key and every log
line in the rest of this chapter.

\subsection{Class hierarchy: MultiNest \emph{is} a Dynesty configuration}
\label{sec:multinest-classes}

\begin{quote}
\sffamily\footnotesize
\code{FeedbackSampler} $\;\to\;$ \code{DynestySampler} $\;\to\;$ \code{MultiNestSampler}
\end{quote}

\code{FeedbackSampler} is the shared base every feedback-driven method in \Jtwo\
builds on (the same one \sampler{AdaptiveBridson} uses, Chapter~\ref{ch:adaptive-levelset})
--- it owns checkpointing and the notion of a synchronisation barrier.
\code{DynestySampler} adds the nested-sampling engine itself, with one class
attribute, \code{default\_dynamic}, that picks the dynamic engine by default.
\code{MultiNestSampler} \emph{is a \code{DynestySampler} subclass that simply flips
that one attribute to \code{False}} and otherwise changes nothing.

\begin{jnote}
\textbf{\sampler{MultiNest} here is not Fortran MultiNest or \code{pymultinest}.}
V1 already used the name \sampler{MultiNest} for a thin, static-only wrapper around
the same engine \sampler{Dynesty} uses --- V2 keeps that exact contract rather than
introduce a second, incompatible meaning for a name your existing cards already use.
If you are starting fresh with no V1 history to preserve, either name is the same
engine; \sampler{MultiNest} simply cannot enable the dynamic engine, ever, even by
accident (\S\ref{sec:nested-validation}).
\end{jnote}

\subsection{The vendored engine and its one patch}
\label{sec:multinest-vendored}

The \code{dynesty} package lives \emph{in-tree}, not as a normal dependency, under
\file{Sampling/Source/Dynesty/}. The only local change is a small patch that threads
a \textsc{uuid} alongside every proposed point through \code{dynesty}'s internal
state (\code{live\_v} / \code{live\_uid}). \code{dynesty} itself has no notion of
``this point was evaluated by Worker~7'' --- vendoring plus this one patch is what
lets \Jtwo\ recover \emph{which} proposed point a returned likelihood value belongs
to, without forking the algorithm itself.

\subsection{The bridge to the Workers: \code{RedisEvaluationPool}}
\label{sec:multinest-pool}

\code{dynesty} is written to be parallel-friendly through one extension point: give
its constructor any object with a \code{.map(func, iterable)} method, and it uses that
for every batch of likelihood evaluations --- ordinarily a
\code{multiprocessing.Pool}. \Jtwo\ hands it \code{RedisEvaluationPool} instead, an
object that \emph{looks} like a local pool to \code{dynesty} but actually round-trips
every evaluation through Redis to the real Worker fleet:

\begin{enumerate}
  \item \textbf{Propose.} \code{dynesty} calls the prior transform on a unit-cube
        vector; \Jtwo's transform mints a fresh \textsc{uuid} and appends it
        (\S\ref{sec:multinest-vendored}) --- this step runs locally, since it is pure
        geometry, not physics.
  \item \textbf{Evaluate.} \code{dynesty} calls \code{pool.map()} with a batch of
        proposed points. \code{RedisEvaluationPool} recognises this as a physics
        batch, builds a real \code{Sample} for each \textsc{uuid}-tagged point, and
        pushes it onto the \emph{same} task queue every other sampler uses
        (Section~\ref{sec:redis-schema}) --- full calculator/opera/likelihood
        pipeline, full \file{SAMPLE} / \file{DATABASE} archiving, full per-Sample
        log, exactly as if a fixed-set sampler had proposed the point.
  \item \textbf{Return.} Once every point in the batch has a result, the pool hands
        \code{dynesty} back a list of floats in the order it expects. \code{dynesty}
        never knows a Worker, a Redis queue, or a calculator was involved.
\end{enumerate}

\begin{jwarn}
\code{dynesty} reuses \code{.map()} for more than likelihood batches --- prior
transforms and some internal proposal bookkeeping go through the same interface
without tagging which is which. An earlier version of the dispatch guess could fail
silently and evaluate your \emph{real} physics function inside the control process,
bypassing every Worker and every calculator. That path now fails loudly instead ---
if the pool cannot prove a batch is a genuine likelihood evaluation, it raises rather
than guesses, so a nested-sampling run either uses your full declared pipeline or
stops immediately, never a plausible-looking wrong answer.
\end{jwarn}

\subsection{From YAML to the engine: pass-through, not reinvention}
\label{sec:multinest-yaml-bridge}

\Jtwo\ does not define its own configuration schema for the engine --- it takes
almost the entire official \code{dynesty} constructor and \code{run\_nested} keyword
surface \emph{as-is} and passes it straight through. Two small filtering functions do
this: one collects constructor keywords (\yamlkey{bound}, \yamlkey{sample},
\yamlkey{walks}, \ldots) from \yamlkey{Bounds} and its nested \yamlkey{sampler} block;
the other collects \yamlkey{run\_nested} keywords, resolving the static/dynamic
evidence-key difference along the way (\S\ref{sec:nested-run-nested}). An unrecognised
key is dropped with a logged warning, never a crash --- and a handful of runtime
handles (\code{loglikelihood}, \code{prior\_transform}, \code{pool}, \code{rstate},
\code{ndim}) are always injected by \Jtwo\ itself and stripped if you paste them into
your card by mistake. This is why the \textsc{yaml} keys in
\S\ref{sec:multinest-yaml} read exactly like the official \code{dynesty}
documentation --- they \emph{are} the official surface.

\subsection{There is no dynamic switch --- Method \emph{is} the switch}
\label{sec:multinest-no-switch}

Earlier builds let you write \yamlkey{Bounds.dynamic: true/false} to pick the engine
under either method name. That flag is gone. \sampler{MultiNest} and
\sampler{Dynesty} share one implementation class (\S\ref{sec:multinest-classes}), and
each subclass now hard-codes which engine it runs; \yamlkey{Sampling.Method} is the
\emph{only} place the choice is made. If a card copied from a \sampler{Dynesty}
example still carries \yamlkey{Bounds.dynamic} after you switch
\yamlkey{Method: MultiNest}, the card is rejected before anything runs
(\S\ref{sec:nested-validation}) rather than silently doing the wrong thing.

\subsection{Running it}
\label{sec:multinest-run-loop}

Put together, one call --- \code{run\_adaptive()} --- does the whole run: build the
\code{RedisEvaluationPool}, seed a \code{numpy} random generator from
\yamlkey{rseed}, construct the (always static, for \sampler{MultiNest})
\code{NestedSampler}, call its \code{run\_nested(**kwargs)}, then export the results
\textsc{csv} and diagnostics \textsc{json}, log the closing summary, and checkpoint
once at that final barrier. That last point matters operationally ---
Section~\ref{sec:nested-distributed} explains why there is no checkpoint
\emph{during} the run.

\section{Configuring MultiNest}
\label{sec:multinest-yaml}

\subsection{Run one, step by step}
\label{sec:nested-run}

\begin{enumerate}
  \item \textbf{Copy the starter template.} \Jtwo\ ships ready-to-edit
        \yamlkey{Sampling} blocks under \code{project\_template/bin/sampling/} in the
        \code{jarvishep2} source tree. For a first run, copy
        \file{Sampling\_MultiNest\_Simple.yaml} into your task card.
  \item \textbf{Replace \yamlkey{Variables} with your parameters}, following the
        distribution catalog of Chapter~\ref{ch:variables}.
  \item \textbf{Replace \yamlkey{LogLikelihood}} with your real expressions
        (Chapter~\ref{ch:loglikelihood}), or wire in calculators/operas
        (Part~\ref{part:backends}) that produce the observables your likelihood needs.
  \item \textbf{Set \yamlkey{Bounds.nlive}} --- the live-point count. Start at
        100--500 for a first look; production evidence estimates typically want
        500--2000.
  \item \textbf{Run it:}
\begin{terminal}
\begin{jcmd}
Jarvis run scan.yaml
\end{jcmd}
\end{terminal}
\end{enumerate}

The minimal working card:

\begin{yamlwin}
Sampling:
  Method: MultiNest
  Variables:
    - name: x
      distribution: {type: Flat, parameters: {min: 0.0, max: 1.0}}
  Bounds:
    Seed: 21
    nlive: 100
    dlogz: 0.5
    run_nested:
      print_progress: true
  LogLikelihood:
    - {name: LogL, expression: "..."}   # your physics
\end{yamlwin}

\subsection{The keys you set on every run}
\label{sec:nested-bounds-meta}

Table~\ref{tab:nested-meta} lists the keys you will set on essentially every
nested-sampling card, whichever engine you chose.

\begin{table}[ht]
  \footnotesize
  \begingroup
  \setlength{\tabcolsep}{0pt}
  \begin{tabular}{@{}p{0.255\linewidth}p{0.213\linewidth}p{0.532\linewidth}@{}}
    \toprule
    \textbf{Key} & \textbf{Default} & \textbf{Meaning} \\
    \midrule
    \yamlkey{nlive} & \code{100} & live-point count $N$ (\S\ref{sec:nested-algorithm});
      the main knob on run cost and posterior resolution (minimum 2) \\
    \yamlkey{Seed} (aliases \yamlkey{seed}, \yamlkey{rseed}) & \code{0} & \textsc{rng}
      seed; fix it for a reproducible proposal sequence. Like every sampler setting it
      goes in \yamlkey{Bounds}, not at the \yamlkey{Sampling} top level \\
    \yamlkey{dlogz} & \code{0.5} & evidence-change stopping threshold
      (\S\ref{sec:nested-dlogz}) \\
    \bottomrule
  \end{tabular}
  \endgroup
  \caption{\yamlkey{Bounds} meta keys, shared by both engines. There is no
  \yamlkey{dynamic} key --- the engine is fixed by \yamlkey{Sampling.Method}
  (\S\ref{sec:multinest-no-switch}); Chapter~\ref{ch:dynesty} covers
  \sampler{Dynesty}'s extra keys.}
  \label{tab:nested-meta}
\end{table}

\begin{jtip}
Tune \yamlkey{nlive} and \yamlkey{dlogz} first, and only those two, until you have a
result. Everything below is for shaping proposals once you know the run is otherwise
healthy.
\end{jtip}

\subsection{Shaping the ellipsoidal bound: \yamlkey{Bounds.sampler}}
\label{sec:nested-sampler-block}

\begin{yamlwin}
Bounds:
  sampler:
    bound: multi      # none | single | multi | balls | cubes
    sample: auto       # auto | unif | rwalk | slice | rslice | hslice
    walks: 25           # rwalk step count
    facc: 0.5             # rwalk target acceptance
    slices: 5               # slice-family step count
    enlarge: 1.25             # bounding-region safety margin
    bootstrap: 0                # bootstrap rounds for the enlargement factor
\end{yamlwin}

\begin{keylist}
  \item[bound] The region proposals are drawn from
        (\S\ref{sec:multinest-ellipsoids}). \code{multi} (the default) is Feroz \&
        Hobson's union of ellipsoids and the right first choice, multimodal or not;
        \code{single} is faster for a simple unimodal posterior; \code{none} samples
        directly from the whole prior (only sensible in very low dimensions).
  \item[sample] How a point is drawn inside that region. \code{auto} (the default)
        picks a sensible method from your dimensionality; \code{unif} is the
        original \sampler{MultiNest} rejection sampler; \code{rwalk} and the
        \code{slice} family scale better to higher dimensions.
\end{keylist}

These are official \code{dynesty} constructor options passed straight through
(\S\ref{sec:multinest-yaml-bridge}); the flat spelling
(\yamlkey{Bounds.bound: multi}) and the nested \yamlkey{Bounds.sampler} block are
both accepted --- the nested block is preferred for readability. The full constructor
surface (periodic/reflective dimensions, custom update intervals, \ldots) is
documented key-by-key in \file{Sampling\_MultiNest\_Full.yaml} beside the simple
template.

\subsection{Controlling the run itself: \yamlkey{Bounds.run\_nested}}
\label{sec:nested-run-nested}

\begin{yamlwin}
Bounds:
  run_nested:
    dlogz: 0.5               # evidence-change stop (static engine)
    add_live: true             # fold surviving live points into Z + posterior
    print_progress: true         # progress lines to the console/log
    maxcall: 100000                 # hard cap on likelihood evaluations
    maxiter: 50000                    # hard cap on nested iterations
\end{yamlwin}

\begin{keylist}
  \item[dlogz] The stopping threshold (\S\ref{sec:nested-algorithm},
        \S\ref{sec:nested-dlogz}). \sampler{MultiNest}'s static engine reads this key
        under the name \yamlkey{dlogz}; \sampler{Dynesty}'s dynamic engine reads the
        \emph{same} \yamlkey{Bounds.dlogz} but internally as
        \yamlkey{dlogz\_init} --- Chapter~\ref{ch:dynesty} explains why.
  \item[add\_live] \opt, default \code{true}. Whether the live points still standing
        at the stop are folded into the evidence and posterior --- almost always what
        you want; the algorithm's own accounting (\S\ref{sec:nested-algorithm}) is
        incomplete without it.
  \item[print\_progress] \opt, default \code{true}. Prints the progress line
        described in \S\ref{sec:nested-log-progress} as the run advances.
  \item[maxcall / maxiter] \opt, no default (unlimited). Hard stops in addition to the
        evidence-based stop --- useful as a safety net on a compute budget.
\end{keylist}

\subsection{Choosing \yamlkey{dlogz}}
\label{sec:nested-dlogz}

\yamlkey{dlogz} sets how precisely the evidence must be pinned down before the run
stops --- smaller values mean a longer run and a tighter uncertainty on $\log Z$.
Table~\ref{tab:dlogz-guide} suggests a value for each kind of run.

\begin{table}[ht]
  \footnotesize
  \begingroup
  \setlength{\tabcolsep}{0pt}
  \begin{tabular}{@{}p{0.34\linewidth}p{0.66\linewidth}@{}}
    \toprule
    \textbf{\yamlkey{dlogz}} & \textbf{Use for} \\
    \midrule
    \code{1.0}   & a quick look --- posterior shape, rough evidence \\
    \code{0.5}   & the default; fine for most exploratory work \\
    \code{0.1}   & a publication-grade evidence estimate \\
    \code{0.01}  & a tight evidence comparison between close models --- expect a much
      longer run \\
    \bottomrule
  \end{tabular}
  \endgroup
  \caption{Choosing \yamlkey{dlogz}. Halving it roughly doubles the run's likelihood
  evaluations for a typical posterior.}
  \label{tab:dlogz-guide}
\end{table}

\section{Validating your card}
\label{sec:nested-validation}

Before \cli{Jarvis run} or \cli{Jarvis check} ever touches Redis, the loaded card
passes through a validation gate that rejects a malformed \yamlkey{Bounds} block
immediately, with a precise, machine-readable diagnosis --- not a stack trace three
minutes into a scan. A clean card prints one line and moves on:

\begin{terminal}
\begin{jout}
Config validation successful. The task YAML meets the V2 contract.
\end{jout}
\end{terminal}

A rejected one lists every problem found, each with a stable code, the exact
\textsc{yaml} path, and (usually) a hint:

\begin{terminal}
\begin{jout}
Config validation failed (1 error(s), 0 warning(s)):
  [error] JV2-BND-012  Sampling.Bounds.dynamic
          Bounds.dynamic is not a V2 setting; the engine is fixed by
          Sampling.Method (Dynesty -> always DynamicNestedSampler,
          MultiNest -> always static NestedSampler). Remove Bounds.dynamic.
          hint: Use Method: Dynesty for dynamic nested sampling, or
          Method: MultiNest for static NestedSampler -- not a dynamic flag.
\end{jout}
\end{terminal}

Check a card without starting anything --- no Redis, no Workers, exit code \code{0}
on success and \code{2} on any error:

\begin{terminal}
\begin{jcmd}
Jarvis validate scan.yaml
Jarvis validate scan.yaml --strict   # also fail on warnings
Jarvis validate scan.yaml --json     # machine-readable report for tooling
\end{jcmd}
\end{terminal}

Table~\ref{tab:nested-validation-codes} lists the findings a nested-sampling card is
most likely to trigger, with the fix for each.

\begin{table}[ht]
  \footnotesize
  \begingroup
  \setlength{\tabcolsep}{0pt}
  \begin{tabular}{@{}p{0.255\linewidth}p{0.213\linewidth}p{0.532\linewidth}@{}}
    \toprule
    \textbf{Code} & \textbf{Fires on} & \textbf{Fix} \\
    \midrule
    \code{JV2-BND-000} & \yamlkey{Bounds} is not a mapping & fix the \textsc{yaml}
      structure \\
    \code{JV2-BND-001} & unknown key directly under \yamlkey{Bounds} & typo, or move
      the key under \yamlkey{sampler}/\yamlkey{run\_nested} \\
    \code{JV2-BND-010} & \yamlkey{nlive} missing or $< 2$ & set an integer $\geq 2$ \\
    \code{JV2-BND-012} & \yamlkey{Bounds.dynamic} / \yamlkey{Dynamic} present at all
      & delete it; switch \yamlkey{Method} instead
      (\S\ref{sec:multinest-no-switch}) \\
    \code{JV2-BND-013} & \yamlkey{dlogz} / \yamlkey{dlogz\_init} not a number &
      fix the value \\
    \code{JV2-BND-021} & unknown key under \yamlkey{run\_nested} & check the static
      vs.\ dynamic key tables (\S\ref{sec:nested-run-nested},
      Chapter~\ref{ch:dynesty}) \\
    \code{JV2-BND-031} & a HEP-injected key (\code{pool}, \code{rstate}, \ldots)
      pasted into \yamlkey{sampler} & remove it --- \Jtwo\ supplies these
      (\S\ref{sec:multinest-yaml-bridge}) \\
    \code{JV2-BND-032} & unknown key under \yamlkey{sampler} / \yamlkey{constructor} &
      check the official \code{dynesty} constructor surface \\
    \code{JV2-VAR-002} & \yamlkey{distribution.type} misspelled or wrong case & names
      are case-sensitive --- \code{Flat}, not \code{flat}
      (Chapter~\ref{ch:variables}) \\
    \code{JV2-MTH-002} & \yamlkey{Sampling.Method} missing on a scan run & set it \\
    \code{JV2-MTH-003} & \yamlkey{Sampling.Method} set to something \Jtwo\ does not
      implement & check spelling against the method list in Section~\ref{sec:method} \\
    \bottomrule
  \end{tabular}
  \endgroup
  \caption{Validation codes a nested-sampling card is most likely to hit. The other
  methods' codes, for reference, are in Appendix~\ref{app:troubleshooting}.}
  \label{tab:nested-validation-codes}
\end{table}

\begin{jnote}
Every error is a hard stop --- there is no partial validation pass for
\yamlkey{Bounds}. Warnings (dead keys such as \yamlkey{Runtime.Subprocess}, an
unrecognised \yamlkey{Operas} \yamlkey{call\_mode}) still let the card run; pass
\cli{--strict} to \cli{Jarvis validate} to treat those as errors too before you
commit a card to a long production run.
\end{jnote}

\section{How it runs on the distributed Workers}
\label{sec:nested-distributed}

Nested sampling behaves differently from every other method in this book, and it is
worth knowing why before you tune \yamlkey{workers}.

\begin{jnote}
The nested-sampling \emph{algorithm} --- which point to retire, where to propose the
next one --- runs entirely in the \textbf{control process}
(\S\ref{sec:multinest-pool}). Only the \textbf{likelihood evaluation} for each
proposed point is sent out to the Workers. The control process waits for that answer
before it can decide its next move.
\end{jnote}

Practically: \yamlkey{EnvReqs.V2.workers} still controls how many likelihood
evaluations run in parallel, and \yamlkey{EnvReqs.V2.batch\_size} still governs how
many are pipelined at once (it becomes the engine's internal \yamlkey{queue\_size}
unless you override it under \yamlkey{Bounds.sampler.queue\_size}). But because the
algorithm itself is single-threaded control logic, adding Workers speeds up the
\emph{evaluation} of each batch of candidates, not the number of \emph{sequential
decisions} nested sampling has to make --- so returns on extra Workers taper off
sooner than for a fixed-set scan of the same size.

\begin{jwarn}
\textbf{Interrupting a nested run does not resume it the way other methods do.} \Jtwo's
own checkpoint/resume system (Chapter~\ref{ch:checkpoint}) records the sampler's state
only \emph{after} the run finishes, for bookkeeping --- there is no mid-run checkpoint
of the live-point population today (\S\ref{sec:multinest-run-loop}). If you need to
survive an interruption on a long run, opt into \code{dynesty}'s own native checkpoint
file:
\begin{yamlwin}
Bounds:
  run_nested:
    checkpoint_file: "&J/outputs/checkpoints/multinest.save"
    checkpoint_every: 60      # seconds
    # resume: true            # set this on the *next* launch to continue
\end{yamlwin}
Leave \yamlkey{resume} unset (or \code{false}) on the first launch; set it to
\code{true} on a later launch pointing at the same \yamlkey{checkpoint\_file} to pick
up where the previous run left off.
\end{jwarn}

\section{Reading the log}
\label{sec:nested-log}

Nested sampling talks more than any other sampler --- deliberately, since a run can
take hours and you want to know it is healthy without waiting for it to finish. Every
excerpt below uses the real Jarvis-logger visual format
(Section~\ref{sec:log-visual-format}); if that format looks unfamiliar, read
Chapter~\ref{ch:logging} first. Everything here lands in \file{sampler.log}
(\emph{not} \file{core.log}).

\subsection{At startup}
\label{sec:nested-log-startup}

Two entries confirm what actually launched, before anything else happens:

\begin{terminal}[sampler.log]
\begin{jout}
·•· Jarvis-HEP.Sampler.MultiNest
	-> 07-21 15:47:05.001 - [INFO] >>>
Starting NestedSampler nlive=100 ndim=2 dynamic=False bound=multi sample=auto

·•· Jarvis-HEP.Sampler.MultiNest
	-> 07-21 15:47:05.004 - [INFO] >>>
run_nested kwargs -> {'dlogz': 0.5, 'add_live': True, 'print_progress': True,
                       'save_bounds': True}
\end{jout}
\end{terminal}

\begin{jtip}
Check this line before you do anything else. \yamlkey{nlive} and \yamlkey{ndim}
should match what you wrote in the card; \yamlkey{dynamic} is not something you set
--- it is \Jtwo\ reporting back which engine \yamlkey{Method} selected, and should
read \code{False} here (\code{True} would mean this is actually a
\sampler{Dynesty} run). The \code{run\_nested kwargs} line is the \emph{final},
resolved configuration --- including any silent defaults
(\S\ref{sec:multinest-yaml-bridge}). If a key you set is missing from this line, it
was not recognised --- look one line up for a matching \code{WARNING} naming it.
\end{jtip}

\subsection{While it runs: the progress block}
\label{sec:nested-log-progress}

With \yamlkey{print\_progress: true} (the default), the \emph{inner} engine logger
--- one dotted segment deeper, \code{...MultiNest.Inner} --- prints a block on every
iteration:

\begin{terminal}[sampler.log]
\begin{jout}
·•· Jarvis-HEP.Sampler.MultiNest.Inner
	-> 07-21 15:47:12.884 - [INFO] >>>
MultiNest Progress ->
	iter       -> 4213
	bound      -> 12
	nc         -> 1
	ncall      -> 5917
	eff(%)     -> 71.204
	loglstar   -> -inf < -3.812 <  inf
	logz       ->  -5.204 +/- 0.183
	dlogz      ->   0.812 >   0.500
\end{jout}
\end{terminal}

Table~\ref{tab:nested-progress-fields} reads that block field by field.

\begin{table}[ht]
  \footnotesize
  \begingroup
  \setlength{\tabcolsep}{0pt}
  \begin{tabular}{@{}p{0.34\linewidth}p{0.66\linewidth}@{}}
    \toprule
    \textbf{Field} & \textbf{Meaning} \\
    \midrule
    \code{iter}      & the current nested iteration --- one retired point per
      iteration (\S\ref{sec:nested-algorithm}) \\
    \code{bound}     & how many times the ellipsoid bound has been rebuilt so far \\
    \code{nc}        & likelihood calls used for the \emph{current} step \\
    \code{ncall}     & total likelihood calls so far --- your actual compute spent \\
    \code{eff(\%)}   & efficiency: \code{iter}/\code{ncall}; healthy runs stay in the
      tens of percent (very low efficiency means the bound is struggling --- see
      below) \\
    \code{loglstar}  & the current worst live point's log-likelihood, between the
      declared bounds \\
    \code{logz}      & the running evidence estimate and its uncertainty --- watch
      this settle over the run \\
    \code{dlogz}     & remaining estimated evidence change, compared against your
      threshold; the run stops once this drops below it \\
    \bottomrule
  \end{tabular}
  \endgroup
  \caption{Progress-block fields, in the order they print.}
  \label{tab:nested-progress-fields}
\end{table}

\begin{jnote}
Every $100^{\text{th}}$ iteration, this exact block is logged at \code{[WARNING]}
instead of \code{[INFO]} --- purely so long-running progress stays visible in a busy
\code{sampler.log} without you having to scroll; it is not an error
(Section~\ref{sec:log-visual-format}).
\end{jnote}

\begin{jtip}
Three numbers tell you almost everything about a run in progress: \code{ncall} rising
steadily means Workers are being fed; \code{logz} settling to a stable value (rather
than still drifting) means the evidence estimate is converging; \code{dlogz}
shrinking toward its threshold tells you how much longer to expect.
\end{jtip}

\begin{jwarn}
\textbf{A stuck \code{ncall} with a moving \code{bound} counter and low \code{eff(\%)}}
usually means the ellipsoid bound is a poor fit for this likelihood surface --- a very
thin or curved high-likelihood ridge, for instance. Try \yamlkey{sample: rslice} or
\yamlkey{sample: hslice} before assuming something is broken; these cost more per
proposal but handle awkward geometries far better than uniform rejection.
\end{jwarn}

\subsection{At the end: the summary}
\label{sec:nested-log-summary}

When the run finishes, one multi-line record and one closing line appear:

\begin{terminal}[sampler.log]
\begin{jout}
·•· Jarvis-HEP.Sampler.MultiNest
	-> 07-21 15:52:40.881 - [INFO] >>>
MultiNest run summary
Summary
=======
nlive: 100
niter: 4213
ncall: 5917
eff(%): 71.204
logz: -5.204 +/- 0.183

·•· Jarvis-HEP.Sampler.MultiNest
	-> 07-21 15:52:40.884 - [INFO] >>>
MultiNest finished logZ=-5.2040 ± 0.1830 niter=4213 ncall=5917
\end{jout}
\end{terminal}

This is your run's headline result --- \code{logz} \textbf{is} the Bayesian evidence
you came for, with its uncertainty. \code{niter} is the number of retired (dead)
points; \code{ncall} is the total likelihood evaluations spent to get there --- larger
than \code{niter} whenever efficiency is below 100\%, since some proposals are
rejected before one lands inside the bound.

\section{What the run leaves behind}
\label{sec:nested-outputs}

Alongside the normal output tree (Chapter~\ref{ch:outputs}), a nested run writes two
extra files under \file{DATABASE/}, listed in Table~\ref{tab:nested-outputs}:

\begin{table}[ht]
  \footnotesize
  \begingroup
  \setlength{\tabcolsep}{0pt}
  \begin{tabular}{@{}p{0.255\linewidth}p{0.213\linewidth}p{0.532\linewidth}@{}}
    \toprule
    \textbf{File} & \textbf{Rows} & \textbf{Contains} \\
    \midrule
    \code{multinest\_result.csv} (\code{dynesty\_result.csv} for
      \sampler{Dynesty}, Chapter~\ref{ch:dynesty}) & one per \textbf{retired point}
      (\code{niter}) & the weighted posterior: \code{uuid}, \code{log\_weight},
      \code{log\_Like}, \code{log\_PriorVolume}, the running \code{log\_Evidence}
      $\pm$ error, and each parameter's physical value \\
    \code{sampler\_summary.json} & one record & the same headline numbers as the
      closing log line --- \code{logz}, \code{logzerr}, \code{ncall}, \code{niter},
      \code{nlive} --- in a form your scripts can parse \\
    \bottomrule
  \end{tabular}
  \endgroup
  \caption{Nested-sampling-specific outputs, in addition to
  \code{DATABASE/samples.hdf5} / \code{samples.csv}.}
  \label{tab:nested-outputs}
\end{table}

\begin{jnote}
\textbf{Two different row counts, on purpose.} \file{samples.hdf5} /
\file{samples.csv} record \emph{every} likelihood evaluation the Workers ever ran
(\code{ncall} rows --- the full archive, Chapter~\ref{ch:outputs}).
\code{multinest\_result.csv} records only the \emph{retired, weighted} points
(\code{niter} rows) --- the actual posterior. For the evidence and the posterior, read
\code{multinest\_result.csv}; for a raw audit of every point ever computed, read
\file{samples.csv}.
\end{jnote}

A stock \JPlot\ scene (\code{dynesty\_runplot}) reads either result \textsc{csv}
directly and draws the standard nested-sampling diagnostic panels --- evidence
trajectory, live-point count, and posterior weight --- with no configuration
(Chapter~\ref{ch:jarvisplot}).

\section{Feedback payload: \yamlkey{FeedbackReturn}}
\label{sec:nested-feedback-return}

Every Worker result for a nested run is boiled down to the one number the algorithm
actually needs before it can decide its next move: the log-likelihood.

\begin{yamlwin}
Sampling:
  FeedbackReturn:
    mode: minimal        # default; correct for MultiNest/Dynesty
    include_logl: true
\end{yamlwin}

The default (\yamlkey{mode: minimal}, \yamlkey{include\_logl: true}) is already what
nested sampling needs, which is why the shipped templates and example projects omit
this block entirely --- write it only if you need a non-default
\yamlkey{FeedbackReturn} for some other reason. A point whose likelihood could not be
computed is fed back as $\log L = -\infty$, which the algorithm treats as an ordinary
rejected proposal rather than a scan failure.

\begin{jtip}
If \code{logz} refuses to settle no matter how long the run goes, and the progress
block's \code{eff(\%)} is healthy, check the likelihood itself first --- a likelihood
that is occasionally $-\infty$ from a failing calculator, rather than from genuinely
being outside your model's support, slows convergence far more than any
\yamlkey{Bounds} tuning can fix.
\end{jtip}

\begin{jtip}
Ready for the dynamic engine? Chapter~\ref{ch:dynesty} assumes everything in this
chapter and covers only what changes.
\end{jtip}

\section{Summary}
\label{sec:multinest-summary}

\sampler{MultiNest} is nested sampling's static engine: a fixed live-point population
that shrinks through the likelihood until its own evidence estimate has converged ---
the only method here that hands you $\log Z$ alongside the posterior, at the cost of no
mid-run checkpoint and \textsc{uuid}s that are not seed-deterministic.
Table~\ref{tab:multinest-summary} gathers it.

\begin{table}[ht]
  \footnotesize
  \begingroup
  \setlength{\tabcolsep}{0pt}
  \begin{tabular}{@{}p{0.34\linewidth}p{0.66\linewidth}@{}}
    \toprule
    \textbf{Field} & \textbf{\sampler{MultiNest}} \\
    \midrule
    Kind & feedback (evidence-stop) \\
    Engine & always static --- not Fortran MultiNest \\
    Point count & until \yamlkey{dlogz} converges, or \yamlkey{maxcall}/\yamlkey{maxiter} \\
    Dimensions & any \\
    Required keys & none strictly; set \yamlkey{Bounds.nlive} \\
    u$\to$x mapping & yes \\
    \yamlkey{Bounds.Seed} role & \code{numpy} generator seed (proposals only) \\
    \textsc{uuid}s & fresh per proposal, \textbf{not} seed-deterministic \\
    Mid-run resume & no --- use \code{dynesty}'s own \yamlkey{checkpoint\_file} \\
    Extra output & \file{multinest\_result.csv}, \file{sampler\_summary.json} \\
    \bottomrule
  \end{tabular}
  \endgroup
  \caption{\sampler{MultiNest} at a glance.}
  \label{tab:multinest-summary}
\end{table}
